\section{Hàm số liên tục}

\subsection{Đề 1}
\Opensolutionfile{ans}[ans/ans-1-C3B3-De1-LC]
\caulc
\begin{ex}%[1D3H3-4]
Cho hàm số $f(x)=\heva{&\dfrac{\sqrt{x}-2}{x-4}&\text{khi}\,\,x\ne 4\\&\dfrac{1}{4}&\text{khi}\,\,x=4}$. Khẳng định nào sau đây \textbf{đúng}?
\choice
{Hàm số có tập xác định $\mathscr{D}=\mathbb{R}\setminus\{4\}$}
{Hàm số liên tục tại mọi điểm trên tập xác định nhưng gián đoạn tại $x=4$}
{Hàm số không liên tục tại $x=4$}
{\True Hàm số liên tục tại $x=4$}
\loigiai{
Tập xác định $\mathscr{D}=[0;+\infty)$.\\
Ta có $\heva{&f(4)=\dfrac{1}{4}\\&\lim\limits_{x\to 4} f(x)=\lim\limits_{x\to 4} \dfrac{1}{\sqrt{x}+2}=\dfrac{1}{4}.}$\\
Vậy hàm số $f(x)$ liên tục tại $x=4$.
}
\end{ex}

\begin{ex}%[1D3H3-3]
Tìm giá trị thực của tham số $m$ để hàm số $f(x)=\heva{&\dfrac{x^3-x^2+2x-2}{x-1}&\text{khi}\,\,x\ne 1\\&3x+m&\text{khi}\,\,x=1\\}$ liên tục tại $x=1$.
\choice
{\True $m=0$}
{$m=2$}
{$m=4$}
{$m=6$}
\loigiai{
Tập xác định $\mathscr{D}=\mathbb{R}$.\\
Theo giả thiết thì ta phải có 
\begin{eqnarray*}
	& & 3+m=f(1)= \lim\limits_{x\to 1} f(x)\\
	&\Leftrightarrow & 3+m=\lim\limits_{x\to 1}\dfrac{x^3-x^2+2x-2}{x-1}\\
	&\Leftrightarrow & 3+m=\lim\limits_{x\to 1} \dfrac{\left(x-1\right)\left(x^2+2 \right) }{x-1}\\
	&\Leftrightarrow &
	3+m=\lim\limits_{x\to 1} \left(x^2+2 \right)=3\\
	&\Leftrightarrow&
	m=0\\
\end{eqnarray*}
}
\end{ex}

\begin{ex}%[1D3H3-3]
Biết rằng hàm số $y=f(x)=\heva{&\dfrac{3-x}{\sqrt{x+1}-2}&\text{khi}\,\,x\ne 3\\&m&\text{khi}\,\,x=3\\}$ liên tục tại $x=3$ (với $m$ là tham số). Khẳng định nào dưới đây \textbf{đúng}?
\choice
{$m\in (-3;0)$}
{\True $m\le -3$}
{$m\in [0;5)$}
{$m\in [0;+\infty)$}
\loigiai{
Tập xác định $\mathscr{D}=\left[-1;+\infty\right)$.\\
Theo giả thiết thì ta phải có 
\begin{eqnarray*}
& & m=f(3)= \lim\limits_{x\to 3} f(x)\\
&\Leftrightarrow & m=\lim\limits_{x\to 3}\dfrac{3-x}{\sqrt{x+1}-2}\\
&\Leftrightarrow &
m=\lim\limits_{x\to 3} \dfrac{\left(3-x\right)\left(\sqrt{x+1}+2\right)}{x-3}\\
&\Leftrightarrow&
m=\lim\limits_{x\to 3}\left[- \left(\sqrt{x+1}+2\right)\right]=-4 \\
\end{eqnarray*}
}
\end{ex}

\begin{ex}%[1D3V3-3]
Số điểm gián đoạn của hàm số $f(x)=\heva{&0{,}5&\text{khi}&\,\,x=-1\\
&\dfrac{x(x+1)}{x^2-1}&\text{khi}&\,\,x\ne-1, x\ne 1\\
&1&\text{khi}&\,\,x=1}$ là
\choice
{$0$}
{\True $1$}
{$2$}
{$3$}
\loigiai{
Tập xác định $\mathscr{D}=\mathbb{R}$.\\
Hàm số $f(x)$ liên tục trên các khoảng $(-\infty;-1)$, $(-1;1)$, $(1;+\infty)$.
\begin{enumerate}[(i)]
\item Xét tại $x=-1$, ta có
\begin{eqnarray*}
\lim\limits_{x\to -1} f(x)
&=&
\lim\limits_{x\to -1} \dfrac{x(x+1)}{x^2-1}\\
&=&
\lim\limits_{x\to -1} \dfrac{x(x+1)}{(x-1)(x+1)}\\
&=&
\lim\limits_{x\to -1}\dfrac{x}{x-1}=\dfrac{1}{2}=f(-1)
\end{eqnarray*}
Vậy hàm số $y=f(x)$ liên tục tại $x=-1$.
\item Xét tại $x=1$, ta có\\
$\heva{&\lim\limits_{x\to 1^{+}} f(x)=\lim\limits_{x\to 1^{+}}\dfrac{x(x+1)}{x^2-1}=\lim\limits_{x\to 1^{+}}\dfrac{x}{x-1}=+\infty\\&\lim\limits_{x\to 1^{-}} f(x)=\lim\limits_{x\to 1^{-}}\dfrac{x(x+1)}{x^2-1}=\lim\limits_{x\to 1^{-}}\dfrac{x}{x-1}=-\infty}$\\
Vậy hàm số $y=f(x)$ gián đoạn tại $x=1$.
\end{enumerate}
}
\end{ex}

\begin{ex}%[1D3V3-4]
Có bao nhiêu giá trị của tham số $a$ để hàm số $f(x)=\heva{&\dfrac{x^2-3x+2}{\left|x-1\right|}&\text{khi}&\,\,x\ne 1\\&a&\text{khi}&\,\,x=1}$ liên tục trên $\mathbb{R}$?
\choice
{$1$}
{$2$}
{\True $0$}
{$3$}
\loigiai{
Tập xác định $\mathscr{D}=\mathbb{R}$.\\
Hàm số $f(x)$ liên tục trên các khoảng $(-\infty;1)$ và $(1;+\infty)$.\\ Khi đó hàm số đã cho liên tục trên $\mathbb{R}$  khi và chỉ khi nó liên tục tại $x=1$, tức là ta cần có $\lim\limits_{x\to 1}f(x)=f(1)\Leftrightarrow \lim\limits_{x\to 1^{+}}f(x)=\lim\limits_{x\to 1^{-}}f(x)=f(1)\quad (\mathrm{I})$\\
Ta có $f(x)=\heva{&x-2&\text{khi}\,\,&x>1\\&a&\text{khi}\,\,&x=1\\&2-x&\text{khi}\,\,&x<1}$\\
Nên $\heva{&\lim\limits_{x\to 1^{-}}f(x)=\lim\limits_{x\to 1^{-}}(2-x)=1\\&\lim\limits_{x\to 1^{+}}f(x)=\lim\limits_{x\to 1^{+}}(x-2)=-1}\Rightarrow (\mathrm{I})$ không thỏa mãn với mọi $x\in \mathbb{R}$.\\
Không tồn tại giá trị $a$ thỏa yêu cầu.
}
\end{ex}

\begin{ex}%[1D3H3-4]
Xét tính liên tục của hàm số $f(x)=\heva{&\dfrac{x-1}{\sqrt{2-x}-1}&\text{khi}&\,\,x<1\\&-2x&\text{khi}&\,\,x\ge 1}$. Khẳng định nào dưới đây \textbf{đúng}?
\choice
{$f(x)$ không liên tục trên khoảng $(-\infty;1)$}
{$f(x)$ không liên tục trên khoảng $(1;+\infty)$}
{$f(x)$ gián đoạn tại $x=1$}
{\True $f(x)$ liên tục trên $\mathbb{R}$}
\loigiai{
Tập xác định $\mathscr{D}=(-\infty;2]$.\\
Xét $x\in (-\infty;1)$, khi đó hàm số $f(x)=\dfrac{x-1}{\sqrt{2-x}-1}$ xác định và liên tục trên khoảng $(-\infty;1)$.\\
Xét $x\in (1;+\infty)$, khi đó hàm số $f(x)=-2x$ xác định và liên tục trên $(1;+\infty)$.\\
Ta có $\heva{&f(1)=-2\\&\lim\limits_{x\to 1^{+}}f(x)=\lim\limits_{x\to 1^{+}}(-2x)=-2\\&\lim\limits_{x\to 1^{-}}f(x)=\lim\limits_{x\to 1^{-}}\dfrac{x-1}{\sqrt{2-x}-1}=\lim\limits_{x\to 1^{-}}\left[-\left(\sqrt{2-x}+1 \right)\right]=-2.}$\\
Vậy hàm số $f(x)$ liên trục trên $\mathbb{R}$.
}
\end{ex}

\begin{ex}%[1D3V3-3]
Tìm giá trị lớn nhất của $a$  để hàm số $f(x)=\heva{&\dfrac{\sqrt[3]{3x+2}-2}{x-2}&\text{khi}&\,\,x>2\\&a^2x-\dfrac{7}{4}&\text{khi}&\,\,x\le 2}$  liên tục tại $x=2$.
\choice
{$a_{\max}=3$}
{$a_{\max}=0$}
{\True $a_{\max}=1$}
{$a_{\max}=2$}
\loigiai{
Tập xác định $\mathscr{D}=\mathbb{R}$.\\
Ta có $\heva{&f(2)=2a^2-\dfrac{7}{4}\\&\lim\limits_{x\to 2^{+}}f(x)=\lim\limits_{x\to 2^{+}}\dfrac{\sqrt[3]{3x+2}-2}{x-2}=\lim\limits_{x\to 2^{+}}\dfrac{3}{\sqrt[3]{(3x+2)^2}+2\sqrt[3]{(3x+2)}+4}=\dfrac{1}{4}\\&\lim\limits_{x\to 2^{-}}f(x)=\lim\limits_{x\to 2^{-}}\left(a^2x+\dfrac{1}{4}\right)=2a^2-\dfrac{7}{4}.}$\\
Hàm số liên tục tại $x=2\Leftrightarrow\lim\limits_{x\to 2^{+}}f(x)=\lim\limits_{x\to 2^{-}}f(x)=f(2)\Leftrightarrow 2a^2-\dfrac{7}{4}=\dfrac{1}{4}\Leftrightarrow a=\pm 1$\\
Vậy $a_{\max}=1$.
}
\end{ex}

\begin{ex}%[1D3V3-4]
Cho hàm số $f(x)=\heva{&\cos\dfrac{\pi x}{2}&\text{khi}&\,\,\left|x\right|\le 1\\&x-1&\text{khi}&\left|x\right| >1 }$.  Mệnh đề nào sau đây là \textbf{sai}?
\choice
{\True Hàm số liên tục tại $x=-1$}
{Hàm số liên tục trên các khoảng $(-\infty;-1)$; $(1;+\infty)$}
{Hàm số liên tục tại $x=1$}
{Hàm số liên tục trên khoảng $(-1;1)$}
\loigiai{
Tập xác định $\mathscr{D}=\mathbb{R}$.\\
Ta có $f(x)$ liên tục trên các khoảng $(-\infty;-1)$, $(-1;1)$, $(1;+\infty)$.\\
Mặt khác $\heva{&f(-1)=\cos\left(-\dfrac{\pi}{2}\right)=0\\&\lim\limits_{x\to (-1)^{+}}f(x)=\lim\limits_{x\to (-1)^{+}}(x-1)=-2.}$ \\
Vậy hàm số $f(x)$ không liên tục tại $x=-1$.
}
\end{ex}

\begin{ex}%[1D3H3-4]
Cho hàm số $f(x)=\heva{&x&\text{khi}&\,\,x<1,x\ne 0\\&0&\text{khi}&\,\,x=0\\&\sqrt{x}&\text{khi}&\,\,x\ge 1}$.  Hàm số $f(x)$ liên tục tại:
\choice
{\True mọi điểm thuộc $\mathbb{R}$}
{mọi điểm trừ $x=1$}
{mọi điểm trừ $x=0$}
{mọi điểm trừ $x=0$ và $x=1$}
\loigiai{
Tập xác định $\mathscr{D}=\mathbb{R}$.\\
Hàm số liên tục trên mỗi khoảng $(-\infty;0)$, $(0;1)$ và $(1;+\infty)$.
\begin{enumerate}[$\bullet$]
\item Tại $x=0$, ta có $\heva{&f(0)=0\\&\lim\limits_{x\to 0}f(x)=\lim\limits_{x\to 0} x=0.}$\\
Vậy hàm số $f(x)$ liên tục tại $x=0$.
\item Tại $x=1$, ta có $\heva{&f(1)=1\\&\lim\limits_{x\to 1^{-}}f(x)=\lim\limits_{x\to 1^{-}}x=1\\&\lim\limits_{x\to 1^{+}}f(x)=\lim\limits_{x\to 1^{+}}\sqrt{x}=1.}$\\
Vậy hàm số $f(x)$ liên tục tại $x=1$.
\end{enumerate}
}
\end{ex}

\begin{ex}%[1D3V3-3]
Số điểm gián đoạn của hàm số $f(x)=\heva{&2x&\text{khi}&\,\,x<0\\&x^2+1&\text{khi}&\,\,0\le x\le 2\\&3x-1&\text{khi}&\,\,x>2}$ là
\choice
{\True $1$}
{$2$}
{$3$}
{$0$}
\loigiai{
Tập xác định $\mathscr{D}=\mathbb{R}$.\\
Hàm số liên tục trên mỗi khoảng $(-\infty;0)$, $(0;2)$ và $(2;+\infty)$.
\begin{enumerate}[$\bullet$]
\item Tại $x=0$, ta có $\heva{&f(0)=1\\&\lim\limits_{x\to 0^{-}}f(x)=\lim\limits_{x\to 0^{-}}2x=0.}$\\
Vậy hàm số $f(x)$ gián đoạn tại $x=0$.
\item Tại $x=2$, ta có $\heva{&f(2)=5\\&\lim\limits_{x\to 2^{-}}f(x)=\lim\limits_{x\to 2^{-}}(x^2+1)=5\\&\lim\limits_{x\to 2^{+}}f(x)=\lim\limits_{x\to 2^{+}}(3x-1)=5.}$\\
Vậy hàm số $f(x)$ liên tục tại $x=2$.
\end{enumerate}
}
\end{ex}

\begin{ex}%[1D3H3-4]
Cho hàm số $f(x)=\heva{&\dfrac{x^2-1}{x-1}&\text{khi}&\,\,x<3,x\ne 1\\&4&\text{khi}&\,\,x=1\\&\sqrt{x+1}&\text{khi}&\,\,x\ge 3}$ . Hàm số $f(x)$ liên tục tại:
\choice
{mọi điểm thuộc $\mathbb{R}$}
{mọi điểm trừ $x=1$}
{mọi điểm trừ $x=3$}
{\True mọi điểm trừ $x=1$ và $x=3$}
\loigiai{
Tập xác định $\mathscr{D}=\mathbb{R}$.\\
Hàm số liên tục trên mỗi khoảng $(-\infty;1)$, $(1;3)$ và $(3;+\infty)$.
\begin{enumerate}[$\bullet$]
\item Tại $x=1$, ta có $\heva{&f(1)=4\\&\lim\limits_{x\to 1}f(x)=\lim\limits_{x\to 1}\dfrac{x^2-1}{x-1}=\lim\limits_{x\to 1}(x+1)=2.}$\\
Vậy hàm số $f(x)$ gián đoạn tại $x=1$.
\item Tại $x=3$, ta có $\heva{&f(3)=2\\&\lim\limits_{x\to 3^{-}}f(x)=\lim\limits_{x\to 3^{-}}\dfrac{x^2-1}{x-1}=\lim\limits_{x\to 3^{-}}(x+1)=4.}$\\
Vậy hàm số $f(x)$ gián đoạn tại $x=3$.
\end{enumerate}
}
\end{ex}

\begin{ex}%[1D3V3-4]
Cho hàm số $f(x)$ liên tục trên đoạn $[-1;4]$ sao cho $f(-1)=2$, $f(4)=7$. Có thể kết luận gì về số nghiệm của phương trình $f(x)=5$ trên đoạn $[-1;4]$.
\choice
{Vô nghiệm}
{\True Có ít nhất một nghiệm}
{Có đúng một nghiệm}
{Có đúng hai nghiệm}
\loigiai{
Ta có $f(x)=5\Leftrightarrow f(x)-5=0$.\\
Đặt $g(x)=f(x)-5$. Khi đó $\heva{&g(-1)=f(-1)-5=2-5=-3\\&g(4)=f(4)-5=7-5=2.}$ \\
$\Rightarrow g(-1)\cdot g(4)=(-3)\cdot 2=-6<0$.\\
Vậy phương trình $g(x)=0$ có ít nhất một nghiệm thuộc khoảng $(-1;4)$ hay phương trình $f(x)=5$ có ít nhất một nghiệm thuộc khoảng $(-1;4)$.
}
\end{ex}
\Closesolutionfile{ans}
\indapan{6}{ans/ans-1-C3B3-De1-LC}

\cauds
\Opensolutionfile{ans}[ans/ans-1-C3B3-De1-DS]
\begin{ex}%[1D3V3-4]
Cho các hàm số $f(x)=\heva{&\dfrac{x^2-4}{x-2}&\text{khi}&\,\,x\ne 2\\&4{,}5&\text{khi}&\,\,x=2}$ và $g(x)=\dfrac{2}{x-1}$. Khi đó.
\choiceTF
{\True Hàm số $g(x)$ liên tục trên khoảng $(-\infty;1)$}
{\True Giới hạn $\lim\limits_{x\to 2}f(x)=4$}
{Hàm số $f(x)$ liên tục tại điểm $x_0=2$}
{Hàm số $\dfrac{f(x)}{g(x)}$ liên tục tại điểm $x_0=2$}
\loigiai{
\begin{itemchoice}
\itemch 
Xét $x_0\in (-\infty;1)$, ta có $g(x_0)=\lim\limits_{x\to x_0}\dfrac{2}{x-1}=\dfrac{2}{x_0-1}$\\
Vậy hàm số $g(x)$ liên tục trên khoảng $(-\infty;1)$.
\itemch Ta có $\lim\limits_{x\to 2}f(x)=\lim\limits_{x\to 2}\dfrac{x^2-4}{x-2}=\lim\limits_{x\to 2}(x+2)=4$.
\itemch Ta có $\heva{&f(2)=4{,}5\\&\lim\limits_{x\to 2}f(x)=\lim\limits_{x\to 2}\dfrac{x^2-4}{x-2}=\lim\limits_{x\to 2}(x+2)=4.}$\\
Vậy hàm số $f(x)$ gián đoạn tại điểm $x_0=2$.
\itemch Do hàm số $f(x)$ gián đoạn tại điểm $x_0=2$ nên hàm số $\dfrac{f(x)}{g(x)}$ cũng gián đoạn tại điểm $x_0=2$. 	
\end{itemchoice}
}
\end{ex}

\begin{ex}%[1D3H3-4]
Xét tính liên tục của các hàm số sau trên tập xác định của nó
\choiceTF
{\True $f(x)=x^3-x^2+8x$  là hàm số liên tục trên $\mathbb{R}$}
{$f(x)=\dfrac{x^2}{x^2-3x}$  là hàm số liên tục trên khoảng $(-\infty;+\infty)$}
{$f(x)=\dfrac{\sin x+1}{x+1}$  là hàm số liên tục trên các khoảng $(-\infty;0)$ và $(0;+\infty)$}
{\True $f(x)=\sqrt{x-2}$  là hàm số liên tục trên nửa khoảng $[2;+\infty)$}
\loigiai{
\begin{itemchoice}
\itemch Vì $f(x)=x^3-x^2+8x$ là hàm đa thức nên hàm số liên tục trên $\mathbb{R}$.
\itemch Vì $f(x)=\dfrac{x^2}{x^2-3x}$ là hàm phân thức xác định trên các khoảng $(-\infty;0)$, $(0;3)$, $(3;+\infty)$ nên hàm số liên tục trên các khoảng $(-\infty;0)$, $(0;3)$,$(3;+\infty)$.\\
Ta có $\heva{&\lim\limits_{x\to 3^{-}}f(x)=\lim\limits_{x\to 3^{-}}\dfrac{x^2}{x^2-3x}=-\infty\\&\lim\limits_{x\to 3^{+}}f(x)=\lim\limits_{x\to 3^{+}}\dfrac{x^2}{x^2-3x}=+\infty.}$\\
Vậy hàm số $f(x)$ gián đoạn tại điểm $x=3$.
\itemch Tập xác định $\mathscr{D}=(-\infty;-1)\cup (-1;+\infty)$.\\ 
Các hàm số $y=\sin x+1$ và $y=x+1$ đều liên tục trên các khoảng $(-\infty;-1)$,$ (-1;+\infty)$.\\
Nên $f(x)$ liên tục trên các khoảng $(-\infty;-1)$,$ (-1;+\infty)$.
\itemch Hàm số xác định trên nửa khoảng $[2;+\infty)$ nên liên tục trên khoảng $(2;+\infty)$.\\ 
Xét $x_0=2$, ta có $\heva{&f(2)=0\\&\lim\limits_{x\to 2^{+}}f(x)=\lim\limits_{x\to 2^{+}}\sqrt{x-2}=0.}$\\
Vậy hàm số $f(x)$ liên tục tại điểm $x_0=2$.
\end{itemchoice}}
\end{ex}

\begin{ex}%[1D3V3-4]
Cho các hàm số $f(x)=\heva{&\dfrac{\sqrt{4x-7}-1}{x^2-4}&\text{khi}&\,\,x> 2\\&\dfrac{5x-9}{2}&\text{khi}&\,\,x\le 2}$ và $g(x)=\heva{&\dfrac{\sqrt{x+2}-2}{2-x}&\text{khi}&\,\,x> 2\\&\dfrac{1-x}{4}&\text{khi}&\,\,x\le 2}$. Khi đó
\choiceTF
{\True Hàm số $f(x)$ liên tục tại điểm $x_0=2$}
{Hàm số $g(x)$ không liên tục trên khoảng $(2;+\infty)$}
{Giới hạn $\lim\limits_{x\to 2^{+}}g(x)=\dfrac{1}{4}$}
{\True Hàm số $y=\dfrac{f(x)}{g(x)}$ liên tục tại điểm $x_0=2$}
\loigiai{
\begin{itemchoice}
\itemch Ta có Ta có $\heva{&f(2)=\dfrac{1}{2}\\&\lim\limits_{x\to 2^{-}}f(x)=\lim\limits_{x\to 2^{-}}\left(\dfrac{5x-9}{2}\right)=\dfrac{1}{2}\\&\lim\limits_{x\to 2^{+}}f(x)=\lim\limits_{x\to 2^{+}}\dfrac{\sqrt{4x-7}-1}{x^2-4}=\lim\limits_{x\to 2^{+}}\dfrac{4}{(x+2)\left(\sqrt{4x-7}+1 \right) }=\dfrac{1}{2}.}$\\
Vậy hàm số $f(x)$ liên tục tại $x_0=2$.
\itemch 
Xét $x_0\in (2;+\infty)$,
ta có\\ $g(x_0)=\lim\limits_{x\to x_0}g(x)=\lim\limits_{x\to x_0}\dfrac{\sqrt{x+2}-2}{2-x}=\lim\limits_{x\to x_0}\left( -\dfrac{1}{\sqrt{x+2}+2}\right)=-\dfrac{1}{\sqrt{x_0+2}+2}$. \\
Vậy hàm số $g(x)$ liên tục trên khoảng $(2;+\infty)$.	
\itemch Ta có $\lim\limits_{x\to 2^{+}}g(x)=\lim\limits_{x\to 2^{+}}\dfrac{\sqrt{x+2}-2}{2-x}=\lim\limits_{x\to 2^{+}}\left( -\dfrac{1}{\sqrt{x+2}+2}\right) =-\dfrac{1}{4}$.
\itemch 
Ta có $g(2)=\lim\limits_{x\to 2^{-}}g(x)=\lim\limits_{x\to 2^{-}}\dfrac{1-x}{4}=-\dfrac{1}{4}$.\\
Do đó hàm số $g(x)$ liên tục tại $x_0=2$.\\
Do các hàm số $f(x)$ và $g(x)$ đều liên tục tại điểm $x_0=2$ và $g(2)=-\dfrac{1}{4}\ne 0$ nên hàm số $y=\dfrac{f(x)}{g(x)}$ cũng liên tục tại điểm $x_0=2$.	
\end{itemchoice}
}
\end{ex}

\begin{ex}%[1D3V3-4]
Cho hàm số $f(x)=\heva{&\dfrac{1-\sqrt{5x+11}}{2x^2-5x-18}&\text{khi}&\,\,x>-2\\&4-x^2&\text{khi}&\,\,x\le -2}$ và $g(x)=\heva{&\dfrac{x^2-x-6}{x+2}&\text{khi}&\,\,x\ne-2\\&2x+a&\text{khi}&\,\,x=-2}$. Khi đó:
\choiceTF
{\True Hàm số $f(x)$ liên tục trên khoảng $(-\infty;-2)$} 
{Hàm số $f(x)$ liên tục tại điểm $x_0=-2$} 
{Để hàm số $g(x)$ liên tục tại điểm $x_0=-2$ thì $a=1$}
{\True Khi $a=-1$ thì hàm số $y=f(x)\cdot g(x)$ gián đoạn tại điểm $x_0=-2$}
\loigiai{
\begin{itemchoice}
\itemch Xét $x_0\in (-\infty;-2)$, ta có $f(x_0)=\lim\limits_{x\to x_0}f(x)=\lim\limits_{x\to x_0}(4-x^2)=4-x_0^2$.\\
Vậy hàm số $f(x)$ liên tục trên khoảng $(-\infty;-2)$.
\itemch 
Ta có $f(-2)=\lim\limits_{x\to {-2}^{-}}f(x)=\lim\limits_{x\to {-2}^{-}}(4-x^2)=0$.\\
Mặt khác
\begin{eqnarray*}
	\lim\limits_{x\to{-2}^{+}}f(x)&=&\lim\limits_{x\to{-2}^{+}}\dfrac{1-\sqrt{5x+11}}{2x^2-5x-18}\\
	&=&	\lim\limits_{x\to{-2}^{+}}\left[-\dfrac{5}{(2x-9)\left(1+\sqrt{5x+11} \right) } \right]\\
	&=&\dfrac{5}{26}
\end{eqnarray*}	
Vậy hàm số $f(x)$ gián đoạn tại điểm $x_0=-2$.
\itemch Ta có $\heva{&g(-2)=-4+a\\&\lim\limits_{x\to {-2}}g(x)=\lim\limits_{x\to {-2}}\dfrac{x^2-x-6}{x+2}=\lim\limits_{x\to {-2}}(x-3)=-5.}$\\
Để hàm số $g(x)$ liên tục tại điểm $x_0=-2\Leftrightarrow -4+a=-5\Leftrightarrow a=-1$.
\itemch 
Do hàm số $f(x)$ gián đoạn tại $x_0=-2$ nên hàm số $y=f(x)\cdot g(x)$ cũng gián đoạn tại điểm $x_0=-2$.
\end{itemchoice}
}		
\end{ex}
\Closesolutionfile{ans}
\indapan{3}{ans/ans-1-C3B3-De1-DS}

\caukq
\Opensolutionfile{ans}[ans/ans-1-C3B3-De1-KQ]
\begin{ex}%[1D3H3-3]
Tìm giá trị của tham số $m$ để hàm số $f(x)=\heva{&\dfrac{x^2-x-2}{x-2}&\text{nếu}&\,\,x\ne 2\\&m+1&\text{nếu}&\,\,x=2}$ liên tục tại điểm $x_0=2$.
\shortans[]{2}
\loigiai{
Tập xác định $\mathscr{D}=\mathbb{R}$.\\
Ta có $\heva{&f\left(2\right)=m+1\\&\lim\limits_{x\to {2}}f(x)=\lim\limits_{x\to {2}}\dfrac{x^2-x-2}{x-2}=\lim\limits_{x\to {2}}(x+1)=3.}$ \\
Để hàm số $f(x)$ liên tục tại $x_0=2\Leftrightarrow f(2)=\lim\limits_{x\to 2}f(x) \Leftrightarrow m+1=3\Leftrightarrow m=2$.
	}
\end{ex}

\begin{ex}%[1D3H3-4]
Tìm giá trị của tham số $a$ để hàm số $f(x)=\heva{&x^2+2x+2&\text{khi}&\,\,x\ge 0\\&x^2+a&\text{khi}&\,\,x<0}$  liên tục trên $\mathbb{R}$.
\shortans[]{2}
\loigiai{
Tập xác định $\mathscr{D}=\mathbb{R}$.\\
Hàm số $f(x)$ liên tục trên các khoảng $(-\infty;0)$, $(0;+\infty)$.\\
Ta có $\heva{&f(0)=\lim\limits_{x\to 0^{+}}f(x)=2\\&\lim\limits_{x\to 0^{-}}f(x)=\lim\limits_{x\to 0^{-}}(x^2+a)=a.}$\\
Hàm số $f(x)$ liên tục trên $\mathbb{R}$ khi và chỉ khi $f(0)=\lim\limits_{x\to 0^{+}}f(x)=\lim\limits_{x\to 0^{-}}f(x)\Leftrightarrow a=2$.
}
\end{ex}

\begin{ex}%[1D3H3-4]
Tìm giá trị của tham số $m$ để hàm số $f(x)=\heva{&\dfrac{x^2-16}{x-4}&\text{khi}&\,\,x>4\\&mx+1&\text{khi}&\,\,x\le 4}$  liên tục trên $\mathbb{R}$.
\shortans[]{$1{,}75$}
\loigiai{
Tập xác định $\mathscr{D}=\mathbb{R}$.\\
Xét  $f(x)=\dfrac{x^2-16}{x-4}$ là hàm phân thức nên liên tục trên khoảng $x\in (4;+\infty)$.\\
$f(x)=mx+1$ là hàm đa thức nên liên tục trên khoảng $x\in (-\infty;4)$.\\
Vậy $f(x)$ trên $\mathbb{R}$ khi và chỉ khi $f(4)=\lim\limits_{x\to 4^{-}}f(x)=\lim\limits_{x\to 4^{+}}f(x)\quad (1)$\\
Ta có $\heva{&f(4)=\lim\limits_{x\to 4^{-}}f(x)=4m+1\\&\lim\limits_{x\to 4^{+}}f(x)=\lim\limits_{x\to 4^{+}}\dfrac{x^2-16}{x-4}=\lim\limits_{x\to 4^{+}}(x+4)=8.}$\\
Hàm số $f(x)$ liên tục trên $\mathbb{R}$ khi và chỉ khi $f(4)=\lim\limits_{x\to 4^{+}}f(x)=\lim\limits_{x\to 4^{-}}f(x)$\\
$\Leftrightarrow 4m+1=8\Leftrightarrow m=\dfrac{7}{4}=1{,}75$.	
}
\end{ex}

\begin{ex}%[1D3H3-3]
Cho hàm số $f(x)=\heva{&\dfrac{x^3-5x+2}{2x^2-x-6}&\text{khi}&\,\,x\ne 2\\&\dfrac{4mx-1}{3}&\text{khi}&\,\,x=2}$ .
Tìm giá trị của tham số $m$  để hàm số liên tục tại $x_0=2$.  
\shortans[]{$0{,}5$}
\loigiai{
Tập xác định $\mathscr{D}=\mathbb{R}\setminus \left\lbrace -\dfrac{3}{2}\right\rbrace $.\\
Ta có $\heva{&f(2)=\dfrac{8m-1}{3}\\&\lim\limits_{x\to 2}f(x)=\lim\limits_{x\to 2}\dfrac{x^3-5x+2}{2x^2-x-6}=\lim\limits_{x\to 2}\dfrac{(x-2)(x^2+2x-1)}{(x-2)(2x+3)}=\lim\limits_{x\to 2}\dfrac{x^2+2x-1}{2x+3}=1.}$\\
Để hàm số $f(x)$ liên tục tại $x_0=2\Leftrightarrow f(2)=\lim\limits_{x\to 2}f(x)\Leftrightarrow \dfrac{8m-1}{3}=1\Leftrightarrow m=\dfrac{1}{2}=0{,}5$.		
}
\end{ex}

\begin{ex}%[1D3H3-3]
Cho hàm số $f(x)=\heva{&\dfrac{x^3-x^2+2x-2}{x-1}&\text{khi}&\,\,x\ne 1\\&3x+m&\text{khi}&\,\,x=1}$.
Tìm giá trị của tham số $m$  để hàm số liên tục tại $x_0=1$.  
\shortans[]{$0$}
\loigiai{
Tập xác định $\mathscr{D}=\mathbb{R}$.\\
Ta có $\heva{&f(1)=3+m\\&\lim\limits_{x\to 1}f(x)=\lim\limits_{x\to 1}\dfrac{x^3-x^2+2x-2}{x-1}=\lim\limits_{x\to 1}(x^2+2)=3.}$\\
Để hàm số $f(x)$ liên tục tại $x_0=1\Leftrightarrow f(1)=\lim\limits_{x\to 1}f(x)\Leftrightarrow 3+m=3\Leftrightarrow m=0$.		
}
\end{ex}

\begin{ex}%[1D3H3-4]
Tìm giá trị của tham số $m$ để hàm số $f(x)=\heva{&x^2+x&\text{khi}&\,\,x<1\\&2&\text{khi}&\,\,x=1\\&mx+1&\text{khi}&\,\,x>1}$ liên tục trên $\mathbb{R}$.
\shortans[]{$1$}
\loigiai{
Tập xác định $\mathscr{D}=\mathbb{R}$.\\
Hàm số $f(x)$ liên tục trên các khoảng $(-\infty;1)$, $(1;+\infty)$.\\ 
Ta có $\heva{&f(1)=2\\&\lim\limits_{x\to 1^{-}}f(x)=\lim\limits_{x\to 1^{-}}(x^2+x)=2\\&\lim\limits_{x\to 1^{+}}f(x)=\lim\limits_{x\to 1^{+}}(mx+1)=m+1.}$\\
Hàm số $f(x)$ liên tục trên $\mathbb{R}$ khi và chỉ khi $f(1)=\lim\limits_{x\to 1^{-}}f(x)=\lim\limits_{x\to 1^{+}}f(x)$.\\
$\Leftrightarrow m+1=2\Leftrightarrow m=1$.
}
\end{ex}
\Closesolutionfile{ans}
\indapan{6}{ans/ans-1-C3B3-De1-KQ}

\subsection{Đề 2}
\Opensolutionfile{ans}[ans/ans-1-C3B3-De2-LC]
\caulc
\begin{ex}%[1D3H3-3]
Cho hàm số $f(x)=\dfrac{x^2-1}{x+1}$ thỏa mãn $f(-1)=m^2-5$. Giá trị của tham số $m$ để hàm số $f(x)$ liên tục tại $x=-1$ là
\choice
{$\sqrt{3}$}
{$-\sqrt{3}$}
{\True $\pm\sqrt{3}$}
{$\pm3$}
\loigiai{
Tập xác định $\mathscr{D}=\mathbb{R}\setminus\{-1\}$.\\
Để hàm số $f(x)$ liên tục tại điểm $x=-1$
\begin{eqnarray*}
&\Leftrightarrow& f(-1)=\lim\limits_{x\to -1} f(x)\\
&\Leftrightarrow& m^2-5=\lim\limits_{x\to -1} \dfrac{x^2-1}{x+1}\\
&\Leftrightarrow&
m^2-5=\lim\limits_{x\to -1}(x-1)\\
&\Leftrightarrow&
m^2-5=-2\\
&\Leftrightarrow&
m=\pm \sqrt{3}
\end{eqnarray*}
}
\end{ex}

\begin{ex}%[1D3H3-3]
	Tìm giá trị thực của tham số $m$ để hàm số $f(x)=\heva{&\dfrac{x^2-x-2}{x-2}&\text{khi}\,\,x\ne 2\\&m&\text{khi}\,\,x=2}$ liên tục tại $x=2$.
	\choice
	{$m=0$}
	{$m=1$}
	{$m=2$}
	{\True $m=3$}
	\loigiai{
Tập xác định $\mathscr{D}=\mathbb{R}$.\\
Hàm số $f(x)$ liên tục trên các khoảng $(-\infty;2)$, $(2;+\infty)$.\\
Theo giả thiết thì ta phải có 
\begin{eqnarray*}
& &m=f(2)= \lim\limits_{x\to 2} f(x)\\	&\Leftrightarrow & m=\lim\limits_{x\to 2}\dfrac{x^2-x-2}{x-2}\\
&\Leftrightarrow & m=\lim\limits_{x\to 2} \left(x+1 \right)\\
&\Leftrightarrow &
m=3
\end{eqnarray*}
	}
\end{ex}

\begin{ex}%[1D3H3-3]
Tìm giá trị thực của tham số $m$ để hàm số $y=f(x)=\heva{&\dfrac{\sqrt{x}-1}{x-1}&\text{khi}\,\,x\ne 1\\&k+1&\text{khi}\,\,x=1\\}$ liên tục tại $x=1$.
\choice
{$k=\dfrac{1}{2}$}
{$k=2$}
{\True $k=-\dfrac{1}{2}$}
{$k=0$}
\loigiai{
Tập xác định $\mathscr{D}=\left[0;+\infty\right) $.\\
Theo giả thiết thì ta phải có 
\begin{eqnarray*}
& & k+1=f(1)= \lim\limits_{x\to 1} f(x)\\
&\Leftrightarrow & k+1=\lim\limits_{x\to 1}\dfrac{\sqrt{x}-1}{x-1}\\
&\Leftrightarrow &
k+1=\lim\limits_{x\to 1} \dfrac{1}{\sqrt{x}+1}=\dfrac{1}{2}\\
&\Leftrightarrow&
k=-\dfrac{1}{2}\\
\end{eqnarray*}
}
\end{ex}

\begin{ex}%[1D3V3-4]
Hàm số $y=f(x)=\heva{&3&\text{khi}&\,\,x=-1\\
&\dfrac{x^4+x}{x^2+x}&\text{khi}&\,\,x\ne-1, x\ne 0\\
&1&\text{khi}&\,\,x=0}$ liên tục tại:
\choice
{mọi điểm trừ $x=0$, $x=1$}
{\True mọi điểm $x\in \mathbb{R}$}
{mọi điểm trừ $x=-1$}
{mọi điểm trừ $x=0$}
\loigiai{
Tập xác định $\mathscr{D}=\mathbb{R}$.\\
Ta có hàm số liên tục trên mỗi khoảng $(-\infty;-1)$, $(-1;0)$ và $(0;+\infty)$.
\begin{enumerate}[(i)]
\item Xét tại $x=-1$, ta có
\begin{eqnarray*}
\lim\limits_{x\to -1} f(x)
&=&
\lim\limits_{x\to -1} \dfrac{x^4+x}{x^2+x}\\
&=&
\lim\limits_{x\to -1} \dfrac{x(x+1)(x^2-x+1)}{x(x+1)}\\
&=&
\lim\limits_{x\to -1}(x^2-x+1)=3=f(-1)
\end{eqnarray*}
Vậy hàm số $y=f(x)$ liên tục tại $x=-1$.
\item Xét tại $x=0$, ta có
\begin{eqnarray*}
\lim\limits_{x\to 0} f(x)
&=&
\lim\limits_{x\to 0} \dfrac{x^4+x}{x^2+x}\\
&=&
\lim\limits_{x\to 0} \dfrac{x(x+1)(x^2-x+1)}{x(x+1)}\\
&=&
\lim\limits_{x\to 0}(x^2-x+1)=1=f(0)
\end{eqnarray*}
Vậy hàm số $y=f(x)$ liên tục tại $x=0$.
\end{enumerate}
}
\end{ex}

\begin{ex}%[1D3V3-4]
Có bao nhiêu giá trị thực của tham số $m$ để hàm số $f(x)=\heva{&m^2x^2&\text{khi}&\,\,x\le 2\\&(1-m)x&\text{khi}&\,\, x>2}$ liên tục trên $\mathbb{R}$?
\choice
{\True $2$}
{$1$}
{$0$}
{$3$}
\loigiai{
Tập xác định $\mathscr{D}=\mathbb{R}$.\\
Ta có hàm số liên tục trên mỗi khoảng $(-\infty;2)$ và $(2;+\infty)$.\\
Ta có $\heva{&f(2)=4m^2\\&\lim\limits_{x\to 2^{+}} f(x)=\lim\limits_{x\to 2^{+}}[(1-m)x]=2(1-m)\\&\lim\limits_{x\to 2^{-}} f(x)=\lim\limits_{x\to 2^{-}}m^2x^2=4m^2}$\\
Hàm số $f(x)$ liên tục trên $\mathbb{R}$ khi và chỉ khi $f(x)$ liên tục tại $x=2$\\
$\Leftrightarrow\lim\limits_{x\to 2^{+}} f(x)=\lim\limits_{x\to 2^{-}} f(x)=f(2)\Leftrightarrow 4m^2=2(1-m)\Leftrightarrow 2m^2+m-1=0\Leftrightarrow \hoac{&m=-1\\&m=\dfrac{1}{2}.}$
	}	
\end{ex}

\begin{ex}%[1D3V3-4]
Biết rằng hàm số  $f(x)=\heva{&\sqrt{x}&\text{khi}&\,\,x\in [0;4]\\&1+m&\text{khi}&\,\, x\in (4;6]}$ liên tục trên $[0;6]$. Khẳng định nào sau đây \textbf{đúng}?
\choice
{\True $m<2$}
{$2\le m<3$}
{$3<m<5$}
{$m \ge 5$}
\loigiai{
Tập xác định $\mathscr{D}=\mathbb{R}$.\\
Hàm số $f(x)$ liên tục trên mỗi khoảng $(0;4)$ và $(4;6)$.\\ Khi đó hàm số liên tục trên đoạn $[0;6]$ khi và chỉ khi hàm số liên tục tại $x=0$, $x=4$, $x=6$.\\
Tức là ta cần có $\heva{&\lim\limits_{x\to 0^{+}}f(x)=f(0)\\&\lim\limits_{x\to 6^{-}}f(x)=f(6)\\&\lim\limits_{x\to 4^{-}}f(x)=\lim\limits_{x\to 4^{+}}f(x)=f(4).}\quad (\mathrm{I})$
\begin{enumerate}[$\bullet$]
\item $\heva{&\lim\limits_{x\to 0^{+}}f(x)=\lim\limits_{x\to 0^{+}}\sqrt{x}=0\\&f(0)=\sqrt{0}=0}$
\item $\heva{&\lim\limits_{x\to 6^{-}}f(x)=\lim\limits_{x\to 6^{-}}(1+m)=1+m\\&f(6)=1+m}$
\item $\heva{&\lim\limits_{x\to 4^{-}}f(x)=\lim\limits_{x\to 4^{-}}\sqrt{x}=2\\&\lim\limits_{x\to 4^{+}}f(x)=\lim\limits_{x\to 4^{+}}(1+m)=1+m\\&f(4)=2}$
\end{enumerate}
Khi đó $(\mathrm{I})$ trở thành $1+m=2\Leftrightarrow m=1$.
}	
\end{ex}

\begin{ex}%[1D3H3-4]
Biết rằng $f(x)=\heva{&\dfrac{x^2-1}{\sqrt{x}-1}&\text{khi}&\,\,x\ne 1\\&a&\text{khi}&\,\,x=1}$ liên tục trên đoạn $[0;1]$ (với $a$ là tham số). Khẳng định nào dưới đây về giá trị $a$ là \textbf{đúng}?
\choice
{\True $a$ là một số nguyên}
{$a$ là một số vô tỉ}
{$a>5$}
{$a<0$}
\loigiai{
Hàm số xác định và liên tục trên $[0;1)$.\\ 
Ta có $\heva{&f(1)=a\\&\lim\limits_{x\to 1^{-}}f(x)=\lim\limits_{x\to 1^{-}}\dfrac{x^2-1}{\sqrt{x}-1}=\lim\limits_{x\to 1^{-}}\left[(x+1)\left(\sqrt{x}+1 \right)\right]=4.}$\\
Hàm số $f(x)$ liên tục trên $[0;1]$ khi và chỉ khi  $\lim\limits_{x\to 1^{-}}f(x)=f(1)\Leftrightarrow a=4$.
}
\end{ex}

\begin{ex}%[1D3H3-3]
Tìm giá trị nhỏ nhất của $a$  để hàm số $f(x)=\heva{&\dfrac{x^2-5x+6}{\sqrt{4x-3}-x}&\text{khi}&\,\,x>3\\&1-a^2x&\text{khi}&\,\,x\le 3}$  liên tục tại $x=3$.
\choice
{\True $-\dfrac{2}{\sqrt{3}}$}
{$\dfrac{2}{\sqrt{3}}$}
{$-\dfrac{4}{3}$}
{$\dfrac{4}{3}$}
\loigiai{
Tập xác định $\mathscr{D}=\mathbb{R}$.\\
Ta có $\heva{&f(3)=1-3a^2\\&\lim\limits_{x\to 3^{+}}f(x)=\lim\limits_{x\to 3^{+}}\dfrac{x^2-5x+6}{\sqrt{4x-3}-x}=\lim\limits_{x\to 3^{+}}\dfrac{(x-2)\left(\sqrt{4x-3}+x\right)}{1-x}=-3\\&\lim\limits_{x\to 3^{-}}f(x)=\lim\limits_{x\to 3^{-}}\left(1-a^2x\right)=1-3a^2.}$\\
Hàm số liên tục tại $x=3\Leftrightarrow\lim\limits_{x\to 3^{+}}f(x)=\lim\limits_{x\to 3^{-}}f(x)=f(3)\Leftrightarrow 1-3a^2=-3\Leftrightarrow a=\pm \dfrac{2}{\sqrt{3}}$.\\
Vậy $a_{\min}=-\dfrac{2}{\sqrt{3}}$.
}
\end{ex}

\begin{ex}%[1D3H3-4]
Xét tính liên tục của hàm số $f(x)=\heva{&1-\cos x&\text{khi}&\,\,x\le 0\\&\sqrt{x+1}&\text{khi}&\,\,x>0}$. Khẳng định nào sau đây \textbf{đúng}?
\choice
{$f(x)$ liên tục tại $x=0$}
{$f(x)$ liên tục trên $(-\infty;1)$}
{\True $f(x)$ không liên tục trên $\mathbb{R}$}
{$f(x)$ gián đoạn tại $x=1$}
\loigiai{
Tập xác định $\mathscr{D}=\mathbb{R}$.\\
Ta có $f(x)$ liên tục trên $(-\infty;0)$ và $(0;+\infty)$,\\
Mặt khác $\heva{&f(0)=0\\&\lim\limits_{x\to 0^{-}}f(x)=\lim\limits_{x\to 0^{-}}(1-\cos x)=0\\&\lim\limits_{x\to 0^{+}}f(x)=\lim\limits_{x\to 0^{+}}\sqrt{x+1}=1}.$\\
Vậy hàm số $f(x)$ không liên tục tại $x=0$ nên $f(x)$ không liên tục trên $\mathbb{R}$.
	}
\end{ex}

\begin{ex}%[1D3N3-2]
Hàm số $f(x)$ có đồ thị như hình bên dưới không liên tục tại điểm có hoành độ là bao nhiêu?
\begin{center}
\begin{tikzpicture}[scale=1,font=\footnotesize,line join=round, line cap=round,>=stealth]
\draw[->] (-2.1,0)--(5.1,0) node[below left] {$x$};
\draw[->] (0,-1.1)--(0,5.1) node[below left] {$y$};
\draw (0,0) node [below left] {$O$};
\foreach \x/\nx in {1/1,2/2}
\draw[thin] (\x,1pt)--(\x,-1pt) node [below] {$\nx$};
\foreach \y/\ny in {1/1,3/3}
\draw[thin] (1pt,\y)--(-1pt,\y) node [left] {$\ny$};
\draw[dashed,thin](1,0)--(1,3)--(0,3);
\draw[dashed,thin](2,0)--(2,1)--(0,1);
\begin{scope}
\clip (-2,-1) rectangle (5,5);
\draw[thick,samples=200,domain=1:-2,smooth,variable=\x] plot (\x,{2*(\x)^2+1});
\draw[thick,samples=200,domain=5:1,smooth,variable=\x] plot (\x,{(\x)-1});
\end{scope}
\end{tikzpicture}
\end{center}
\choice
{$x=0$}
{\True $x=1$}
{$x=2$}
{$x=3$}
\loigiai{
Dựa vào đồ thị, ta thấy $\heva{&\lim\limits_{x\to 1^{+}}f(x)=0\\&\lim\limits_{x\to 1^{-}}f(x)=3}$ nên hàm số $f(x)$ gián đoạn tại $x=1$.
}
\end{ex}

\begin{ex}%[1D3H3-3]
Tính tổng $S$ gồm tất cả các giá trị $m$ để hàm số $f(x)=\heva{&x^2+x&\text{khi}&\,\,x<1\\&2&\text{khi}&\,\,x=1\\&m^2x+1&\text{khi}&\,\,x>1}$  liên tục tại $x=1$. 
\choice
{$S=-1$}
{\True $S=0$}
{$S=1$}
{$S=2$}
\loigiai{
Tập xác định $\mathscr{D}=\mathbb{R}$.\\
Ta có $\heva{&f(1)=2\\&\lim\limits_{x\to 1^{-}}f(x)=\lim\limits_{x\to 1^{-}}(x^2+x)=2\\&\lim\limits_{x\to 1^{+}}f(x)=\lim\limits_{x\to 1^{+}}(m^2x+1)=m^2+1.}$\\
Hàm số $f(x)$ liên tục tại $x=1$ khi và chỉ khi $\lim\limits_{x\to 1^{+}}f(x)=\lim\limits_{x\to 1^{-}}f(x)=f(1)$\\
$\Leftrightarrow m^2+1=2\Leftrightarrow m=\pm 1$.
}
\end{ex}

\begin{ex}%[1D3C3-5]
Có tất cả bao nhiêu giá trị nguyên của tham số $m$ thuộc khoảng $(-10;10)$ để phương trình $x^3-3x^2+(2m+2)x+m-3=0$  có ba nghiệm phân biệt $x_1,x_2,x_3$  thỏa mãn $x_1<-1<x_2<x_3$?
\choice
{$19$}
{$18$}
{\True $4$}
{$3$}
\loigiai{
Xét hàm số $f(x)=x^3-3x^2+(2m+2)x+m-3$  liên tục trên $\mathbb{R}$.\\
Giả sử phương trình có ba nghiệm phân biệt $x_1,x_2,x_3$ sao cho $x_1<-1<x_2<x_3$. Khi đó $f(x)=(x-x_1)(x-x_2)(x-x_3)$.\\
Ta có $f(-1)=(-1-x_1)(-1-x_2)(-1-x_3)>0$ (do $x_1<-1<x_2<x_3$).\\
Mà $f(-1)=-m-5$ nên suy ra $-m-5>0\Leftrightarrow m<-5$.\\
Thử lại với $m<-5$, ta có \\
$\lim\limits_{x\to -\infty}f(x)=-\infty$ nên tồn tại $a<-1$ sao cho $f(a)<0.\quad (1)$\\
Do $m<-5$ nên $f(-1)=-m-5>0.\quad (2)$\\
Mặt khác $f(0)=m-3<0\quad (3)$\\
$\lim\limits_{x\to +\infty}f(x)=+\infty$ nên tồn tại $b>0$ sao cho $f(b)>0.\quad (4)$\\
Từ $(1)$ và $(2)$, suy ra phương trình có nghiệm thuộc khoảng $(-\infty;-1)$.\\
Từ $(2)$ và $(3)$, suy ra phương trình có nghiệm thuộc khoảng $(-1;0)$.\\
Từ $(3)$ và $(4)$, suy ra phương trình có nghiệm thuộc khoảng $(0;+\infty)$.\\ 
Vậy $m<-5$ thỏa mãn yêu cầu bài toán.\\
Mà $\heva{&-10<m<10\\&m\in \mathbb{Z}}\Rightarrow m\in \left\lbrace -9;-8;-7;-6 \right\rbrace $. 
	}
\end{ex}	
\Closesolutionfile{ans}
\indapan{6}{ans/ans-1-C3B3-De2-LC}

\cauds
\Opensolutionfile{ans}[ans/ans-1-C3B3-De2-DS]
\begin{ex}%[1D3H3-3]
Cho hàm số $f(x)=\heva{&\dfrac{x^2-1}{x-1}&\text{khi}&\,\,x\ne 1\\&x+1&\text{khi}&\,\,x=1}$  và $g(x)=4x^2-x+1$. Khi đó:
\choiceTF
{\True Ta có $f(1)=2$}
{\True Hàm số $f(x)$ liên tục tại điểm $x_0=1$ }
{\True Hàm số $g(x)$ liên tục tại điểm $x_0=2$ }
{Hàm số $y=f(x)-g(x)$ không liên tục tại điểm $x_0=1$}
\loigiai{
\begin{itemchoice}
\itemch Ta có $f(1)=1+1=2$.	
\itemch Ta có $\heva{&f(1)=2\\&\lim\limits_{x\to 1}f(x)=\lim\limits_{x\to 1}\dfrac{x^2-1}{x-1}=\lim\limits_{x\to 1}(x+1)=2.}$\\
Vậy hàm số $f(x)$ liên tục tại điểm $x_0=1$.
\itemch Ta có $\heva{&g(2)=15\\&\lim\limits_{x\to 2}g(x)=\lim\limits_{x\to 2}(4x^2-x+1)=15.}$\\
Vậy hàm số $g(x)$ liên tục tại điểm $x_0=2$.
\itemch 
Ta có $\heva{&g(1)=4\\&\lim\limits_{x\to 1}g(x)=\lim\limits_{x\to 1}(4x^2-x+1)=4.}$\\
Vậy hàm số $g(x)$ liên tục tại điểm $x_0=1$.\\
Do các hàm số $f(x)$ và 	$g(x)$ liên tục tại điểm $x_0=1$ nên hàm số $y=f(x)-g(x)$ cũng liên tục tại điểm $x_0=1$.	
\end{itemchoice}	
}
\end{ex}

\begin{ex}%[1D3V3-3]
Cho các hàm số  $f(x)=\heva{&-\dfrac{x}{2}&\text{khi}&\,\,x\le 1\\&\dfrac{x^2-3x+2}{x^2-1}&\text{khi}&\,\,x>1}$, $g(x)=x^2-3x+1$ và $h(x)=\sin\dfrac{\pi x}{4}$. Khi đó:
\choiceTF
{\True Hàm số $f(x)$ liên tục tại điểm $x_0=1$}
{\True Hàm số $g(x)$ liên tục tại điểm $x_0=2$}
{Hàm số $h(x)$ không liên tục tại điểm $x_0=2$}
{Hàm số $y=f(x)\cdot g(x)$ không liên tục tại điểm $x_0=1$}
\loigiai{
\begin{itemchoice}
\itemch Ta có $\heva{&f(1)=-\dfrac{1}{2}\\&\lim\limits_{x\to 1^{-}}f(x)=\lim\limits_{x\to 1^{-}}\left(-\dfrac{x}{2}\right) =-\dfrac{1}{2}\\&\lim\limits_{x\to 1^{+}}f(x)=\lim\limits_{x\to 1^{+}}\dfrac{x^2-3x+2}{x^2-1}=\lim\limits_{x\to 1^{+}}\dfrac{x-2}{x+1} =-\dfrac{1}{2}.}$\\
Vậy hàm số $f(x)$ liên tục tại $x_0=1$.
\itemch Ta có $\heva{&g(2)=-1\\&\lim\limits_{x\to 2}g(x)=\lim\limits_{x\to 2}\left(x^2-3x+1\right) =-1.}$\\
Vậy hàm số $g(x)$ liên tục tại $x_0=2$.
\itemch Ta có $\heva{&h(2)=1\\&\lim\limits_{x\to 2}h(x)=\lim\limits_{x\to 2}\left(\sin\dfrac{\pi x}{4}\right) =1.}$\\
Vậy hàm số $h(x)$ liên tục tại $x_0=2$.
\itemch 
Ta có $\heva{&g(1)=-1\\&\lim\limits_{x\to 1}g(x)=\lim\limits_{x\to 1}\left(x^2-3x+1\right) =-1.}$\\
Vậy hàm số $g(x)$ liên tục tại $x_0=1$.\\
Do các hàm số $f(x)$ và $g(x)$ đều liên tục tại $x_0=1$ nên hàm số $y=f(x)\cdot g(x)$ cũng liên tục tại $x_0=1$.
\end{itemchoice}
}
\end{ex}

\begin{ex}%[1D3V3-3]
Cho hàm số $f(x)=\heva{&\dfrac{\sqrt{x-1}-1}{x^2-3x+2}&\text{khi}&\,\,x\ne 2\\&\dfrac{2a+1}{6}&\text{khi}&\,\,x=2}$ và $g(x)=\sin \dfrac{\pi x}{4}$. Khi đó:
\choiceTF
{\True Giới hạn $\lim\limits_{x\to 2}f(x)=\dfrac{1}{2}$}
{\True Hàm số $g(x)$ liên tục tại điểm $x_0=2$}
{\True Khi $a=1$ thì hàm số $f(x)$ liên tục tại $x_0=2$}
{Khi $a=0$ thì hàm số $y=f(x)+g(x)$ liên tục tại $x_0=2$}
\loigiai{
\begin{itemchoice}
\itemch  Ta có $\lim\limits_{x\to 2}f(x)=\lim\limits_{x\to 2}\dfrac{\sqrt{x-1}-1}{x^2-3x+2}=\dfrac{1}{(x-1)\left(\sqrt{x-1}+1 \right)}=\dfrac{1}{2}$
\itemch  Ta có $\heva{&g(2)=1\\&\lim\limits_{x\to 2}g(x)=\lim\limits_{x\to 2}\sin \dfrac{\pi x}{4}=1.}$\\
Vậy hàm số $g(x)$ liên tục tại điểm $x_0=2$.
\itemch Khi $a=1$, ta có $\heva{&f(2)=\dfrac{2\cdot1+1}{6}=\dfrac{1}{2}\\&\lim\limits_{x\to 2}f(x)=\lim\limits_{x\to 2}\dfrac{\sqrt{x-1}-1}{x^2-3x+2}=\lim\limits_{x\to 2}\dfrac{1}{(x-1)\left(\sqrt{x-1}+1 \right)}=\dfrac{1}{2}.}$\\
Vậy hàm số $f(x)$ liên tục tại điểm $x_0=2$.
\itemch Khi $a=0$, ta có $\heva{&f(2)=\dfrac{2\cdot0+1}{6}=\dfrac{1}{6}\\&\lim\limits_{x\to 2}f(x)=\lim\limits_{x\to 2}\dfrac{\sqrt{x-1}-1}{x^2-3x+2}=\lim\limits_{x\to 2}\dfrac{1}{(x-1)\left(\sqrt{x-1}+1 \right)}=\dfrac{1}{2}.}$\\
Vậy hàm số $f(x)$ gián đoạn tại điểm $x_0=2$.\\
Do đó hàm số $y=f(x)+g(x)$ cũng gián đoạn tại điểm $x_0=2$.
\end{itemchoice}		
}
\end{ex}

\begin{ex}%[1D3V3-4]
Cho hàm số $f(x)=\heva{&\dfrac{2-\sqrt{x+5}}{x^2-5x-4}&\text{khi}&\,\,x>-1\\&x^2-9x&\text{khi}&\,\,x\le -1}$ và $g(x)=\heva{&\dfrac{x^2-1}{x+1}&\text{khi}&\,\,x\ne-1\\&2a+1&\text{khi}&\,\,x=-1}$. Khi đó:
\choiceTF
{\True Ta có $\lim\limits_{x\to {-1}^{+}}f(x)=\dfrac{1}{8}$} 
{\True Hàm số $f(x)$ liên tục trên khoảng $(-\infty;-1)$} 
{Hàm số $g(x)$ liên tục tại điểm $x_0=-1$ thì $a=\dfrac{1}{2}$}
{ Khi $a=-\dfrac{1}{2}$ thì hàm số $y=f(x)-g(x)$ liên tục tại điểm $x_0=-1$}
\loigiai{
\begin{itemchoice}
\itemch Ta có
\begin{eqnarray*}
\lim\limits_{x\to{-1}^{+}}f(x)&=&\lim\limits_{x\to{-1}^{+}}\dfrac{2-\sqrt{x+5}}{x^2-5x-4}\\
&=&	\lim\limits_{x\to{-1}^{+}}\left[-\dfrac{1}{(x-4)\left(2+\sqrt{x+5} \right) } \right]\\
&=&\dfrac{1}{8}
\end{eqnarray*}	
\itemch 
Xét $x_0\in(-\infty;-1)$, ta có $f(x_0)=\lim\limits_{x\to x_0}f(x)=x_0^2-9x_0$.\\
Vậy hàm số $f(x)$ liên tục trên khoảng $(-\infty;-1)$.
\itemch Ta có $\heva{&g(-1)=2a+1\\&\lim\limits_{x\to {-1}}g(x)=\lim\limits_{x\to {-1}}\dfrac{x^2-1}{x-1}=\lim\limits_{x\to {-1}}(x+1)=0.}$\\
Để hàm số $g(x)$ liên tục tại điểm $x_0=-1\Leftrightarrow 2a+1=0\Leftrightarrow a=-\dfrac{1}{2}$.
\itemch 
Ta có $\heva{&f(-1)=\lim\limits_{x\to {-1}^{-}}f(x)=10\\&\lim\limits_{x\to {-1}^{+}}f(x)=\dfrac{1}{8}}$.\\
Vậy hàm số $f(x)$ gián đoạn tại $x_0=-1$.\\
Do hàm số $f(x)$ gián đoạn tại $x_0=-1$ nên hàm số $y=f(x)-g(x)$ cũng gián đoạn tại điểm $x_0=-1$.
\end{itemchoice}
}		
\end{ex}
\Closesolutionfile{ans}
\indapan{3}{ans/ans-1-C3B3-De2-DS}

\caukq
\Opensolutionfile{ans}[ans/ans-1-C3B3-De2-KQ]
\begin{ex}%[1D3H3-3]
Có bao nhiêu giá trị của tham số $m$ để hàm số $f(x)=\heva{&\dfrac{\sqrt{x}-1}{x^2-1}&\text{nếu}&\,\,x\ne 1\\&m^2x&\text{nếu}&\,\,x=1}$ liên tục tại điểm $x_0=1$?
\shortans[]{2}
\loigiai{
Tập xác định $\mathscr{D}=[0;+\infty)$.\\
Ta có $\heva{&f\left(1\right)=m^2\\&\lim\limits_{x\to 1}f(x)=\lim\limits_{x\to 1}\dfrac{\sqrt{x}-1}{x^2-1}=\lim\limits_{x\to 1}\dfrac{1}{(x+1)\left(\sqrt{x}+1\right) }=\dfrac{1}{4}.}$ \\
Để hàm số $f(x)$ liên tục tại $x_0=1\Leftrightarrow f(1)=\lim\limits_{x\to 1}f(x)\Leftrightarrow m^2=\dfrac{1}{4}\Leftrightarrow m=\pm \dfrac{1}{2}$.
}
\end{ex}

\begin{ex}%[1D3H3-3]
Tìm giá trị của tham số $a$ để hàm số $f(x)=\heva{&\dfrac{\sqrt{x}-1}{x-1}&\text{khi}&\,\,x>1\\&ax-\dfrac{1}{2}&\text{khi}&\,\,x\le 1}$  liên tục tại điểm $x_0=1$.
\shortans[]{1}
\loigiai{
Tập xác định $\mathscr{D}=\mathbb{R}$.\\
Ta có $\heva{&f(1)=\lim\limits_{x\to 1^{-}}f(x)=a-\dfrac{1}{2}\\&\lim\limits_{x\to 1^{+}}f(x)=\lim\limits_{x\to 1^{+}}\dfrac{\sqrt{x}-1}{x-1}=\lim\limits_{x\to 1^{+}}\dfrac{1}{\sqrt{x}+1}=\dfrac{1}{2}.}$\\
Hàm số $f(x)$ liên tục tại điểm $x_0=1\Leftrightarrow f(1)=\lim\limits_{x\to 1^{+}}f(x)=\lim\limits_{x\to 1^{-}}f(x)$\\
$\Leftrightarrow a-\dfrac{1}{2}=\dfrac{1}{2}\Leftrightarrow a=1$.
}
\end{ex}

\begin{ex}%[1D3H3-6]
Một chất điểm chuyển động với tốc độ được cho bởi hàm số\\ $v(t)=\heva{&m+5&\text{khi}&\,\,0\le t\le 5\\&t^2-5t+10&\text{khi}&\,\,t>5}$, 
trong đó $v(t)$ được tính theo đơn vị $\mathrm{m/s}$ và $t$ được tính theo giây ($m$ là tham số). Tìm giá trị của tham số $m$ để hàm số $v(t)$ liên tục tại $t=5$.
\shortans[]{$5$}
\loigiai{
Tập xác định $\mathscr{D}=\mathbb{R}$.\\
Ta có $\heva{&v(5)=\lim\limits_{x\to 5^{-}}v(t)=m+5\\&\lim\limits_{x\to 5^{+}}v(t)=\lim\limits_{x\to 5^{+}}(t^2-5t+10)=10.}$\\
Để hàm số $v(t)$ liên tục tại $t=5\Leftrightarrow v(5)=\lim\limits_{t\to 5}v(t)\Leftrightarrow m+5=10\Leftrightarrow m=5$.	
}
\end{ex}

\begin{ex}%[1D3H3-3]
Cho hàm số $f(x)=\heva{&x^2+x+1&\text{khi}&\,\,x\ne 4\\&2a+1&\text{khi}&\,\,x=4}$ .
Tìm giá trị của tham số $a$  để hàm số liên tục tại $x_0=4$.  
\shortans[]{$10$}
\loigiai{
Tập xác định $\mathscr{D}=\mathbb{R}$.\\
Ta có $\heva{&f(4)=2a+1\\&\lim\limits_{x\to 4}f(x)=\lim\limits_{x\to 4}(x^2+x+1)=21.}$\\
Để hàm số $f(x)$ liên tục tại $x_0=4\Leftrightarrow f(4)=\lim\limits_{x\to 4}f(x)\Leftrightarrow 2a+1=21\Leftrightarrow a=10$.		
}
\end{ex}

\begin{ex}%[1D3H3-4]
Tìm giá trị của tham số $m$ để hàm số $f(x)=\heva{&\dfrac{x^2-x-2}{x-2}&\text{khi}&\,\,x\ne 2\\&m&\text{khi}&\,\,x=2}$ liên tục trên $\mathbb{R}$.
\shortans[]{$3$}
\loigiai{
Tập xác định $\mathscr{D}=\mathbb{R}$.\\
Hàm số $f(x)$ liên tục trên các khoảng $(-\infty;2)$, $(2;+\infty)$.\\
Ta có $\heva{&f(2)=m\\&\lim\limits_{x\to 2}f(x)=\lim\limits_{x\to 2}\dfrac{x^2-x-2}{x-2}=\lim\limits_{x\to 2}(x+1)=3.}$\\
Hàm số $f(x)$ liên tục trên $\mathbb{R}$ khi và chỉ khi $f(2)=\lim\limits_{x\to 2}f(x)\Leftrightarrow m=3$.
}
\end{ex}

\begin{ex}%[1D3H3-3]
Tìm giá trị của tham số $m$ để hàm số $f(x)=\heva{&\dfrac{3-x}{\sqrt{x+1}-2}&\text{khi}&\,\,x\ne 3\\&m&\text{khi}&\,\,x=3}$ liên tục tại $x_0=3$.
\shortans[]{$-4$}
\loigiai{
Tập xác định $\mathscr{D}=[-1;+\infty)$.\\
Ta có $\heva{&f(3)=m\\&\lim\limits_{x\to 3}f(x)=\lim\limits_{x\to 3}\dfrac{3-x}{\sqrt{x+1}-2}=\lim\limits_{x\to 3}\left[ -\left(\sqrt{x+1}+2\right)\right]=-4.}$\\
Để hàm số $f(x)$ liên tục tại điểm $x_0=3\Leftrightarrow f(3)=\lim\limits_{x\to 3}f(x)\Leftrightarrow m=-4$.
}
\end{ex}
\Closesolutionfile{ans}
\indapan{6}{ans/ans-1-C3B3-De2-KQ}
